\documentclass[11pt,a4paper]{article}

\usepackage[margin=2.5cm]{geometry}
\usepackage{fontspec}
\setmainfont{Latin Modern Roman}
\setsansfont{Latin Modern Sans}
\setmonofont{Hack}[Scale=0.85]
\usepackage{amssymb}
\usepackage{amsmath}
\usepackage{graphicx}
\usepackage[dvipsnames,table]{xcolor}
\usepackage{tikz}
\usetikzlibrary{positioning,arrows.meta,fit,backgrounds,shapes.geometric,calc,decorations.pathmorphing}
\usepackage{listings}
\usepackage{enumitem}
\usepackage{hyperref}
\usepackage{fancyhdr}
\usepackage{titlesec}
\usepackage{parskip}
\usepackage{float}
\usepackage{booktabs}
\usepackage{array}
\usepackage{adjustbox}
\usepackage[most]{tcolorbox}

\definecolor{router}{HTML}{E74C3C}
\definecolor{admitted}{HTML}{2ECC71}
\definecolor{denied}{HTML}{E74C3C}
\definecolor{tls}{HTML}{2ECC71}
\definecolor{cipher}{HTML}{9B59B6}
\definecolor{session}{HTML}{3498DB}
\definecolor{ca}{HTML}{F39C12}
\definecolor{codebg}{HTML}{F8F8F8}
\definecolor{codeframe}{HTML}{DDDDDD}
\definecolor{nixblue}{HTML}{5277C3}
\definecolor{soft}{HTML}{F4F6F8}

\newtcolorbox{abstractbox}[1][]{
  enhanced, colback=soft, colframe=nixblue,
  boxrule=0.6pt, arc=2pt, left=10pt, right=10pt, top=8pt, bottom=8pt,
  title=\textbf{Abstract}, coltitle=white, colbacktitle=nixblue,
  fonttitle=\sffamily\bfseries, #1
}
\newtcolorbox{findingbox}[1][]{
  enhanced, colback=soft, colframe=router,
  boxrule=0.6pt, arc=2pt, left=10pt, right=10pt, top=8pt, bottom=8pt,
  title=\textbf{Finding}, coltitle=white, colbacktitle=router,
  fonttitle=\sffamily\bfseries, #1
}

\lstset{
  basicstyle=\ttfamily\small,
  backgroundcolor=\color{codebg},
  frame=single,
  rulecolor=\color{codeframe},
  breaklines=true,
  breakatwhitespace=true,
  tabsize=2,
  showstringspaces=false,
}

\pagestyle{fancy}
\fancyhf{}
\fancyhead[L]{\small\textit{Iterating on the Zenoh security model}}
\fancyhead[R]{\small\textit{DVerse Platform}}
\fancyfoot[C]{\thepage}
\renewcommand{\headrulewidth}{0.4pt}

\hypersetup{
  colorlinks=true,
  linkcolor=blue!60!black,
  urlcolor=blue!60!black,
}

\title{%
  \vspace{-1cm}
  {\LARGE\bfseries Iterating on the Zenoh Security Model:\\
  Three Hypotheses, Two Disconfirmations, One Pivot}
}
\author{Yordan Mitev}
\date{\today}

\begin{document}
\maketitle
\thispagestyle{fancy}

\begin{abstractbox}
  Across three sprints I iterated the Zenoh security model for the
  DVerse platform from open pub/sub, through mTLS authentication,
  to admission-driven Megolm encryption. Each iteration started
  from a hypothesis about what was enough; each of the first two
  was disconfirmed by something specific that the next iteration
  set out to fix. This document is the condensed version of that
  journey: what each step looked like, why it stopped being
  enough, and what replaced it. The full evaluation lives in
  \emph{Zenoh evaluation and security architecture}, and the code
  is described in \emph{Framework architecture and integration}.
\end{abstractbox}

\vspace{1em}

\begin{table}[H]
  \centering
  \small
  \begin{tabular}{|l|l|p{8cm}|}
    \hline
    \rowcolor{gray!15}
    \textbf{Version} & \textbf{Date} & \textbf{Changes} \\
    \hline
    V1 & 9 June 2026 & Initial writeup of the V1 \(\rightarrow\)
    V2 \(\rightarrow\) V3 iteration, the central ``ACLs are not
    enough'' finding, the method and results, and the remaining
    open questions. \\
    \hline
  \end{tabular}
  \caption{Revision history.}
  \label{tab:iter-history}
\end{table}

\newpage
\tableofcontents
\newpage

% ==============================================================
\section{Framing}
% ==============================================================

The DVerse design challenge is a decentralised communication
platform for humans and AI agents, built on Zenoh and Rust. Security
of communication is the part of the challenge I owned. The plan at
sprint 1 was to spend a couple of sprints learning Zenoh's built-in
security model, then commit to a design. The plan held for the
first two sprints. It did not survive contact with sprint 3.

What follows is a short trace of three iterations on the security
model, structured the way each iteration actually happened: a
hypothesis about what was enough, an attempt to ship it, and the
finding that closed it off as a final state.

% ==============================================================
\section{V1: Open Pub/Sub}
% ==============================================================

\textbf{Hypothesis.} A flat Zenoh fabric with a known router
endpoint is enough for a research prototype. The platform is being
prototyped end-to-end and the security model can be slotted in
later.

\textbf{What shipped.} A ping-pong demo in Rust using stock Zenoh.
Two agents, no authentication, plaintext payloads. Documented in
V1 of the security doc.

\textbf{Disconfirmation.} The first non-localhost demo. Anyone on
the same LAN can subscribe to any topic by default, because Zenoh
ships with multicast scouting on by default. The
``decentralised but private'' goal in the design challenge fell
apart at the first network-attached demo with another laptop on
the same Wi-Fi.

V1 was the easy iteration. The disconfirmation was expected; the
purpose of V1 was to learn the Zenoh API and to have a working
ping-pong to layer security onto.

% ==============================================================
\section{V2: mTLS}
% ==============================================================

\textbf{Hypothesis.} Pin every router and every agent to a leaf
certificate issued by a private CA. Configure Zenoh's transport
layer to require client authentication. Only cert holders can
join the fabric, and that solves the V1 leak.

\textbf{What shipped.} A step-ca instance backed by Keycloak (OIDC
identity), short-lived (24-hour) leaf certificates for every
operator and every agent, and a Zenoh transport that rejects
connections without a valid CA-signed cert at the TLS handshake.
Mutual authentication, both ends present a cert. Documented in V2
of the security doc.

\textbf{Disconfirmation.} mTLS authenticates the fabric, but
every cert holder is on the same fabric. A platform that hosts
multiple parallel sessions (one team's working group, another
team's working group, the admin's own session) cannot use mTLS
alone to keep team A out of team B's traffic. They all have valid
certs from the same CA. The fabric is private; the groups inside
it are not.

The fix that V2 sketched in its ``Future Work'' section was Zenoh's
ACL configuration. ACLs can match on the cert's common name and
allow or deny per topic per CN. The plan was to admit a CN to a
session by writing it into the session's ACL block; deny by
removing it. Then I tried to ship that, and ran into the central
finding of this experiment.

% ==============================================================
\section{The Finding: ACLs Are Read Once}
% ==============================================================

\begin{findingbox}
Zenoh's ACL configuration is consumed when the Zenoh session is
opened. There is no API to mutate the admitted-CN list afterwards,
and no published RFC to add one. Changing who can publish or
subscribe to a topic requires building a new \texttt{Config},
closing the running \texttt{Session}, and calling
\texttt{zenoh::open(config)} again. While the session is closed,
no message gets through.
\end{findingbox}

This is the discovery that closed V2 as a final state. For a
system whose admin makes accept/deny decisions at runtime, ACL
mutation is the common path on the data plane. Doing that with
static ACLs means the data plane is interrupted every time a
member joins or leaves. The continuous availability that the
platform exists to provide is violated each time someone uses it.

The router code in \texttt{src/zenoh\_router/src/router.rs} carries
this restriction openly. The outer \texttt{run()} loop returns
from \texttt{session\_loop} whenever the admitted set changes,
aborts the admission handler, closes the Zenoh session, marks
itself \texttt{RouterStatus::Reloading}, and iterates. A new
session opens against the updated ACL. Until it opens, the router
is offline on the wire.

Numerically, the cost is impossible to make small. Even a perfect
session rebuild takes hundreds of milliseconds (cert reload,
listener bind, peer reconnect, subscriber re-declare). With a
realistic admission workload (a handful of joins during a working
session), the result is a noticeable choppiness that nothing in
the application can route around.

The conclusion was that ACLs cannot be the primary access-control
layer of a system with runtime admission. They are workable as a
fabric-level gate that changes rarely; they are not workable as
the per-action knob.

% ==============================================================
\section{V3: Admission Plus Megolm}
% ==============================================================

\textbf{Hypothesis.} Move the per-action access control off the
fabric and into a cryptographic layer above it. Admission becomes
a handshake. Payloads are encrypted with a per-session Megolm
\texttt{SessionKey}. A peer that holds the cert but does not hold
the session key sees opaque ciphertext, even where the ACL would
let it subscribe.

\textbf{What shipped.}

\begin{itemize}[leftmargin=1.4em]
  \item An admission protocol with two well-known topics under
        \path{dverse/session/}: the requester publishes a signed
        \texttt{JoinRequest} on \path{requests/<cn>}, the admin
        replies with an \texttt{Allow} or \texttt{Deny} decision
        on \path{admission/<cn>}.
  \item An ECDSA P-256 binding signature that ties the
        requester's vodozemac Curve25519 identity key to its TLS
        cert's CN. The admin verifies the signature against the
        cert before deciding.
  \item An Olm pre-key wrap of the admin's Megolm
        \texttt{SessionKey}, addressed to the requester's
        one-time prekey. Only the requester can decrypt it; once
        decrypted, the requester installs a
        \texttt{GroupReceiver} keyed by the admin's session id.
  \item Per-payload Megolm AEAD on every agent message. The
        cipher dispatches incoming envelopes by their embedded
        \texttt{sender\_id} to the correct \texttt{GroupReceiver}.
  \item A cross-machine relay on
        \path{dverse/agent_keys/**} that republishes each agent's
        outbound \texttt{SessionKey} every five seconds, with a
        subscriber that installs each envelope idempotently.
  \item A coarsened fabric ACL: one broad allow rule for
        \texttt{dverse/**} scoped to the admitted CNs. The rule
        still rebuilds the Zenoh session on real membership
        changes, but membership changes are now the common case,
        not the per-action case.
\end{itemize}

\textbf{Outcome.} Verified end-to-end on a two-machine LAN with
two operators, each running ping and pong agents. Cross-machine
ping-pong decrypts on the first agent-key relay tick (sub-second
in practice). Adding a new operator to the session adds two
network round-trips (the join request and the admission
decision); no payloads are dropped on the data plane during
admission because the data plane is encrypted, not gated.

% ==============================================================
\section{Method}
% ==============================================================

In rough order:

\begin{enumerate}[leftmargin=1.6em]
  \item Read Zenoh's ACL implementation to confirm the
        finding. Searched for any reload, refresh, or update API.
        There is none. The ACL block is read in
        \texttt{zenoh::open} and not consulted again.
  \item Observed the rebuild cost manually. Running ping-pong
        between two operators, an admin's
        \texttt{admit\_request} drops the agent on the admin side
        for several seconds (the session rebuild) while ping
        keeps producing payloads that never arrive. The
        ``micro-outage'' is loud enough to see in the agent
        logs.
  \item Designed the admission handshake on paper. Picked
        Olm/Megolm because the project's threat model maps
        closely to Matrix's (group communication, per-message
        forward secrecy, key delivery as part of group join).
        Adopted vodozemac as the Rust implementation.
  \item Designed the cryptographic identity binding so that a
        compromise of an unrelated Curve25519 key cannot
        impersonate an operator. Domain-separated, length-prefixed
        message; ECDSA P-256 signature with the TLS leaf's
        private key; verifier extracts the public key from the
        leaf cert.
  \item Implemented the wire format
        (\texttt{src/bot\_framework/src/admission.rs}) and the
        handshake handler
        (\texttt{src/zenoh\_router/src/admission\_handler.rs}).
  \item Implemented the per-process Megolm cipher
        (\texttt{src/bot\_framework/src/payload\_crypto.rs}) and
        the same-machine filesystem mechanism for key sharing.
  \item Found, while testing across two physical machines, that
        same-machine peers decrypt each other and cross-machine
        peers do not. Added the agent-key wire relay to cover
        that case.
  \item Found, while testing logout-then-rejoin flows, that the
        router silently kept its previous session config across
        restarts. Added the \texttt{config\_changed} doorbell to
        force the router to consume the staged config on every
        outer-loop iteration.
\end{enumerate}

% ==============================================================
\section{Results}
% ==============================================================

Table~\ref{tab:iter-result-shape} summarises the cost shape of
each iteration. The numbers are qualitative because the experiment
is observational, not benchmarked; quantitative measurement is a
follow-up if needed.

\begin{table}[H]
  \centering
  \small
  \begin{tabular}{|p{5.4cm}|p{1.4cm}|p{2.6cm}|p{4cm}|}
    \hline
    \rowcolor{gray!15}
    \textbf{Operation} & \textbf{V1} & \textbf{V2} & \textbf{V3} \\
    \hline
    Outsider sees plaintext on wire & yes & no (mTLS) & no (mTLS) \\
    \hline
    Outsider with cert sees plaintext & N/A & yes & no (Megolm) \\
    \hline
    Per-admission data-plane interruption & N/A & yes (session rebuild) & no (key install only) \\
    \hline
    Per-payload crypto cost & none & TLS only & TLS + Megolm AEAD \\
    \hline
    Cross-machine peer-key discovery & N/A & N/A & 5\,s republish; sub-second typical \\
    \hline
  \end{tabular}
  \caption{Cost shape of each iteration.}
  \label{tab:iter-result-shape}
\end{table}

The end-to-end demonstration was a two-machine, two-operator
session. On machine A, operator ``test1'' running an admin router,
a ping agent, and a pong agent. On machine B, operator ``test2''
running a client router, a ping agent, and a pong agent. test1
admitted test2 from the launcher's pending-requests panel; within
five seconds (the agent-key relay tick) every agent on every
machine was decrypting messages from every other agent on every
other machine. No session rebuild was observed during the
admission decision.

% ==============================================================
\section{What Remained Open}
% ==============================================================

V3 settles the core question (``can we have runtime admission
without per-join data-plane outages''), but leaves three items as
follow-on work.

\textbf{Per-topic ACLs as second-line defence.} The fabric rule
admits the broad \texttt{dverse/**} expression for all admitted
CNs. A compromised admitted member can therefore flood any topic;
the Megolm layer keeps the floods unreadable, but they still
consume bandwidth. Adding per-topic ACL rules would let the
fabric drop traffic from CNs that have no business sending it
even though they are otherwise admitted.

\textbf{Signed agent-key envelopes.} The
\texttt{dverse/agent\_keys/**} relay accepts every envelope from
every admitted CN. A compromised admitted member could spam
fabricated session keys, bloating peers' receiver maps. Signing
the envelope with the publishing operator's identity key would
let receivers drop envelopes that fail verification.

\textbf{Multi-session on one fabric.} V3 keeps the historic
``one fabric per session'' shape from V2. A more general system
would allow multiple sessions to share an mTLS fabric and rely on
their Megolm keys for isolation. This works in principle (each
session has its own \texttt{SessionKey}; cross-session encryption
fails by design) but the launcher's chooser-and-state model is
not yet structured for it.

% ==============================================================
\section{Cross-References}
% ==============================================================

The full evaluation of Zenoh's security model, including V1's
ping-pong validation, V2's TLS handshake walkthrough, and the
deep dive into V3's admission machinery, is in
\emph{Zenoh evaluation and security architecture}. The Rust
implementation that backs V3 is described in
\emph{Framework architecture and integration}: which modules in
the \texttt{bot\_framework} crate hold which piece of the
handshake, and how the router, the launcher, and the agents
compose against them.

% ==============================================================
\section*{Glossary}
\addcontentsline{toc}{section}{Glossary}
% ==============================================================

\begin{description}[leftmargin=2.6cm, style=nextline, font=\sffamily\bfseries]
  \item[ACL] Access Control List. In Zenoh, a JSON block in
        the session configuration that defines which subjects
        are allowed or denied on which key expressions. Read
        once at session-open time; not reloadable at runtime.
  \item[admission] The handshake by which a candidate node
        becomes a member of a DVerse session and obtains the
        session's Megolm key.
  \item[AEAD] Authenticated Encryption with Associated Data.
        The cipher mode used by Megolm; provides
        confidentiality and integrity in a single primitive.
  \item[CN] Common Name. The subject field of a TLS
        certificate. The DVerse fabric identifies operators
        by the CN in their leaf certificate.
  \item[Curve25519] An elliptic-curve key-exchange primitive
        used by Olm. The V3 identity binding ties an
        operator's Curve25519 key to the TLS certificate via
        an ECDSA signature.
  \item[fabric] The set of mutually authenticated Zenoh peers
        connected through the same mTLS-protected mesh.
  \item[\texttt{GroupReceiver}] The Megolm decrypt half. One
        is installed per known peer.
  \item[Keycloak] An open-source identity-and-access
        management product; the OIDC provider used by V3
        certificate enrolment.
  \item[Megolm] A group-ratchet cryptographic primitive from
        the Matrix project, used in V3 to encrypt every agent
        payload.
  \item[mTLS] Mutual Transport Layer Security. Both ends of
        the TLS connection authenticate with certificates.
  \item[OIDC] OpenID Connect. The identity layer on top of
        OAuth 2.0 used by Keycloak and Step-CA to validate
        enrolment.
  \item[Olm] A one-to-one cryptographic ratchet from the
        Matrix project, used in V3 to wrap the Megolm
        session key for a specific requester at admission
        time.
  \item[prekey] A short-lived Curve25519 key published by the
        admission requester. Consumed once by the admin's
        Olm session.
  \item[\texttt{SessionKey}] The Megolm key material
        exchanged during admission and installed by the
        requester as a \texttt{GroupReceiver}.
  \item[Step-CA] An online certificate authority from
        Smallstep; configured with an OIDC provisioner that
        validates JWTs against Keycloak.
  \item[TLS] Transport Layer Security. The DVerse fabric
        uses TLS 1.3 with both ends authenticated
        (\emph{mTLS}).
  \item[Zenoh] A decentralised pub/sub messaging protocol
        from the Eclipse Foundation. The DVerse data plane
        is built on Zenoh sessions.
\end{description}

\end{document}
